
% This thesis outlines a series of investigations into the complex environment of the Coma Cluster. The results from these three studies converge on a broader narrative of galaxy and cluster evolution in dense environments. The spatial and luminosity distribution of \glspl{ucd} suggests a composite origin shaped by both internal cluster dynamics and hierarchical assembly. The analysis and identification of galaxies with anomalously low \gls{gc} populations raises key questions about environmental suppression or stripping effects in galaxy clusters that challenge conventional models of \gls{gc} formation. 
% Finally, testing the correlation of metallicity between \glspl{css} and their host galaxies offers a crucial diagnostic of origin, distinguishing between in situ formation and accretion. Through exploration of the manner in which \gls{css} evolve in, and are affected by, dense cluster environments, we ensure these investigations impose meaningful constraints on the processes that govern star cluster populations, galaxy evolution, and the buildup of large-scale structure in rich clusters like Coma.
This thesis uses \glspl{css}---particularly \glspl{gc} and \glspl{ucd}---as long-lived tracers of environmental processing in the dense core of the Coma Cluster. Using deep \gls{hst}/ACS imaging, we analyze the spatial distributions, luminosity functions, color subpopulations, and azimuthal structure of \gls{css} populations across the central $\qty{\sim1}{\mega\parsec}$.

In the first study, we examine the bright end of the \gls{gclf} and the cluster-wide \gls{ucd} distribution. We identify an excess of bright \glspl{css} relative to an extrapolated \gls{gclf} and find that blue \glspl{ucd} are more spatially extended than red \glspl{ucd}, consistent with a composite \gls{ucd} population that may include an environmentally produced contribution (\eg stripped nuclei).

In the second study, we compare expected and observed \gls{gc} populations for major Coma galaxies using two-dimensional Voronoi \gls{gc} density maps. We identify several galaxies with depleted \glspl{gcs} and show that their projected \gls{gcs} morphologies range from strongly asymmetric to nearly isotropic, highlighting degeneracies between recent stripping, older phase mixing, symmetric truncation, and intrinsic formation efficiency.

In the third study, we introduce a blue-\gls{gc} fraction excess diagnostic and relate it to cluster-centric radius, dual-\gls{bcg} influence, and the local \gls{gc} density field. The strongest trends are linked to local \gls{gc} density and \gls{bcg}-related measures, implying that the Coma core is a structured, multi-scale environment rather than a smooth radially isotropic field. 

Overall, the results show that environmental processing in Coma is local and population-dependent, acting most strongly on the extended blue \gls{gc} component and coupling galaxy-halo transformation to the buildup of the intracluster \gls{css} population.

% Include this line is using the glossary and abbreviations - this will reset the 'first-use' of any glossary terms after abstract use.
\glsresetall 